\section{Vad hårdvaran tar över}
\label{sec:handover}

CAN löser samma problem som \kapref{ch:frames}--\kapref{ch:bussar}, i kisel:

\begin{booktable}{@{}lL@{}}
\thead{Mekanism} & \thead{Vem gör den i CAN} \\ \midrule
Ramsynkronisering & Kontrollern: SOF-biten, och sju recessiva bitar som slutmärke. \\
Längdfält & Kontrollern: DLC, fyra bitar, 0--8 databytes. \\
Checksumma & Kontrollern: CRC-15, beräknad och kontrollerad av samma motor. \\
Kollisionshantering & Bussen själv, genom arbitrering --- utan bussmästare och utan förlorad tid. \\
Kvittens & Kontrollern: en enda bit inuti framen, som varje mottagare drar ned. \\
Adressering & \textbf{Finns inte.} Identifieraren beskriver innehållet, och varje nod filtrerar
  själv. \\
Sekvensnummer & \textbf{Finns inte.} Behövs det får applikationen lägga det i nyttolasten. \\
Timeout och omsändning & En fullständig kontroller gör det i hårdvara. \emph{Den här gör det
  inte} --- se \secref{sec:limits}. \\
\end{booktable}

Det som återstår för mjukvaran är alltså att lämna över en identifierare, en längd och några
databytes --- och att läsa av hur det gick.

\begin{aside}
\note CAN som protokoll --- bussen, framen, bitstoppningen, CRC-15, arbitreringen, bittajmingen
och felhanteringen --- står i sin helhet i kursens CAN-bok. Det här kapitlet förutsätter dess
kapitel 1, 2 och 7. Läs dem först; resten av boken är fördjupning du kan ta när den behövs.
\end{aside}

\section{Kontrollern, och varför ett lager ligger emellan}
\label{sec:layer}

Kontrollern är konstruerad i VHDL av en parallellklass. Den är byggd för en anropare i samma
klockdomän, och signalerar därför med pulser som är \textbf{en klockcykel} långa --- 20~ns vid
50~MHz.

En processor som pollar över en seriebuss ser systemet några tiotals mikrosekunder i taget i bästa
fall. En puls på 20~ns är alltså borta tusentals cykler innan nästa pollning kommer.

Därför ligger det två lager mellan kontrollern och din kod:

\begin{console}
                 din C++-kod
                      |
                     SPI                  5-bytetransaktioner, <= 1 MHz
                      |
   spi_slave  ->  spi_reg_bridge          bytes blir transaktioner
                      |
              register_bank               pulser blir klibbiga, pollbara bitar
                      |
              can_controller              CAN i VHDL
                      |
                  CAN-bussen
\end{console}

\textbf{Registerbanken är lagret som gör konstruktionen användbar från mjukvara.} Den fångar
pulserna och håller kvar dem som nivåer du själv måste nollställa. Varje "skriv \code{0x1} här för
att kvittera" i registerkartan är den mekanismen.
